\begingroup\fontsize{9.25}{11.4}\selectfont\setlength{\tabcolsep}{3pt}
\setlength{\tabcolsep}{5pt}
\renewcommand{\arraystretch}{1.18}\begin{longtable}{@{}>{\raggedright\arraybackslash}p{0.23\linewidth}rrrrr@{}}
\caption{Heterogeneity reporting gates across all six outcome families}\label{tab:pub_A26} \\
\toprule Wave / level & IV cells & Pass gate & Adj. / signals & Assignment cells & Adj. / signals \\ \midrule\endfirsthead
\toprule Wave / level & IV cells & Pass gate & Adj. / signals & Assignment cells & Adj. / signals \\ \midrule\endhead
2013 CCPP & 48 & 0 & 0 / -- & 48 & 40 / 2 \\
2013 Household & 48 & 0 & 0 / -- & 48 & 40 / 1 \\
2013 Individual & 64 & 8 & 8 / 0 & 64 & 55 / 3 \\
2017 CCPP & 48 & 24 & 24 / 0 & 48 & 40 / 2 \\
2017 Household & 48 & 8 & 8 / 0 & 48 & 40 / 6 \\
2017 Individual & 64 & 23 & 23 / 0 & 64 & 55 / 1 \\
\bottomrule\end{longtable}{\footnotesize\textit{Notes:} IV counts use the registered primary common-window fuzzy interaction specification. A reportable IV contrast must pass rank, support and conditional-instrument-strength checks, including conditional F statistics above 10. Assignment contrasts instead require supported local estimation; the fuzzy-IV gate does not apply. Adj. counts eligible contrasts with available multiplicity-adjusted probabilities; signals count those below 0.05. The correction is Holm for primary moderators and BH for secondary moderators, within original registered families, not one new pooled family. A dash means no adjusted contrast is available, not a negative finding. Separate rdhte assignment contrasts do not identify treatment-receipt heterogeneity. Sources: all six module-05 aggregate result files.}\par\endgroup
